problem:  
Let \( \mathcal{N}(n) \) be the variety of nilpotent Lie brackets on \( \mathbb{R}^n \), with \( \mathrm{GL}(n) \) acting by change of basis.  In Lauret’s moment–map framework, each bracket \( \mu\in\mathcal{N}(n) \) has a moment map \(m(\mu)\in\operatorname{Sym}(n)\), and the functional \(F(\mu)=\|m(\mu)\|^2\) is \( \mathrm{GL}(n)\)-equivariant up to scale.  Its critical points correspond to **nilsoliton brackets**, which are also **distinguished (minimal) orbits** in geometric invariant theory.  
For every \(\beta\in\operatorname{Sym}(n)\) arising as a normalized mean curvature vector, one obtains a **stratum** \(\mathcal{S}_\beta\subset\mathcal{N}(n)\) containing exactly those brackets whose negative gradient flow of \(F\) converges to an orbit with moment map proportional to \(\beta\).

> **Problem (Nilsoliton Strata Classification Problem).**  
> For a fixed dimension \(n\), determine all strata \(\mathcal{S}_\beta\) that contain nilsoliton orbits, and describe whether a single stratum can contain **more than one** non‑equivalent nilsoliton bracket (modulo scaling and \(\mathrm{O}(n)\)).  
>  
> Equivalently: beyond the known uniqueness of the minimal orbit inside a single orbit closure, is the *moment‑map type* \(\beta\) uniquely determined by the isomorphism class of the nilsoliton, or can distinct nilsoliton Lie algebras correspond to critical points lying in one and the same instability stratum?

---

Why it is a "good" problem:  

1. **Mathematical precision and consistency:**  
   The moment‑map and GIT structure are used correctly here: orbit closures flow to unique minimal points, but the stratification by \(\mathcal{S}_\beta\) may contain several minimal (critical) orbits.  The problem asks precisely about this multiplicity, avoiding contradictions with known results on uniqueness within orbit closures.

2. **Authentic research‑level unknown:**  
   While existence and uniqueness of nilsolitons within a given orbit are established, the higher‑level pattern—how many distinct nilsoliton orbits occupy one stratum, and how \(\beta\) encodes algebraic type—is not completely classified.  Answering this would deepen understanding of the global landscape of Einstein solvmanifolds across dimensions.

3. **Bridge between fields:**  
   The question unites **geometric invariant theory** (stratification by instability type) and **differential geometry** (nilsolitons ↔ homogeneous Ricci solitons ↔ Einstein solvmanifolds).  Studying the distribution of nilsoliton orbits across moment‑map strata connects algebraic stability notions with geometric curvature data.

4. **Simplicity with depth:**  
   The statement is brief—*"Can a single moment‑map stratum contain more than one nilsoliton type?"*—but resolving it demands analysis of the moment–map energy landscape, eigenvalue patterns of Einstein derivations, and orbit‑type geometry.

5. **Potential impact:**  
   - A proof that each stratum carries at most one nilsoliton would imply extremely strong rigidity of homogeneous Einstein moduli.  
   - Conversely, examples of multiple nilsolitons per stratum would reveal hidden degeneracy mechanisms and subtle connections among different solvable Einstein geometries.

Thus the problem is logically sound, aligned with known theory, yet opens a genuine and profound question about the global arrangement of Einstein solvmanifolds within the GIT stratification, exemplifying the depth and clarity of a "good" research‑level problem.